\chapter{Resumes, Interviews, and Career Preparation}

\section{Your Capstone Project as a Career Asset}
Capstone is the richest source of resume content you have as a graduating senior. Unlike coursework, capstone involves real clients, real budgets, real deadlines, interdisciplinary teamwork, and technical problem-solving under uncertainty. Employers recognize this. When they ask about ``relevant project experience'' in an interview, your capstone project is the answer.

\section{Writing a Capstone-Informed Resume}
Quantify your contributions wherever possible. Numbers are more compelling than adjectives:

\begin{itemize}[nosep,leftmargin=1.5em]
\item ``Managed a \$3,500 project budget with procurement of 47 components across 8 vendors, delivering under budget by 12\%.''
\item ``Designed and tested a PID temperature controller achieving $\pm$0.3\,\degC{} accuracy against a $\pm$1.0\,\degC{} requirement.''
\item ``Led a team of 5 engineers through PDR, CDR, and FDR for an external industry client over two semesters.''
\item ``Designed a 4-layer PCB with mixed-signal layout, ordered from JLCPCB, and completed board bring-up with zero rework.''
\end{itemize}

List technical skills demonstrated: specific software (SolidWorks, MATLAB, KiCad, LabVIEW, Python, ANSYS), fabrication methods (CNC machining, 3D printing, PCB assembly, soldering), analysis techniques (FEA, CFD, thermal analysis, statistical testing), and professional skills (design reviews, client presentations, technical writing).

\section{Career Day Preparation}
Career Day falls during SDI (no class that day). This is a prime networking opportunity. Prepare: update your resume (emphasize capstone contributions even if the project is early-stage), practice your elevator pitch, dress professionally (business casual minimum), and bring 10--15 printed resumes.

Use your capstone project as a conversation starter. Employers are interested in your \emph{process} (how you approached the problem, how you made decisions, how you handled setbacks) as much as the \emph{outcome}. Saying ``I am working on a capstone project designing a solar-powered irrigation controller for a community garden, and I am responsible for the sensor integration and PID control'' immediately communicates initiative, technical skill, and teamwork.

\section{Behavioral Interviews and the STAR Method}
Behavioral interview questions ask you to describe past experiences as evidence of your skills and character. Common capstone-relevant questions:

\begin{itemize}[nosep,leftmargin=1.5em]
\item ``Tell me about a time you had a conflict with a teammate. How did you resolve it?''
\item ``Describe a project where something went significantly wrong. How did you handle it?''
\item ``Give an example of when you had to make a technical decision with incomplete information.''
\item ``Tell me about a time you had to deliver bad news to a stakeholder.''
\item ``Describe a situation where you had to learn a new skill quickly to complete a task.''
\end{itemize}

Use the \textbf{STAR method}: \textbf{S}ituation (set the context (your project, the phase, the stakeholders), \textbf{T}ask (what was your specific responsibility), \textbf{A}ction (what did you specifically do) not what the team did), \textbf{R}esult (what was the outcome, and what did you learn). Capstone provides abundant material for STAR responses. Prepare 5--7 STAR stories before interview season and practice telling each in 90 seconds.

\section{Technical Interviews}
Some employers conduct technical interviews where they assess your engineering depth. Your capstone project is excellent material. Be prepared to explain: your system architecture (block diagram level), the key engineering decisions (why did you choose this sensor over that one?), the analysis you performed (show you understand the physics), and the test results (including what did not work and what you learned). Interviewers are assessing whether you think like an engineer, not whether your project was perfect.

\section{Networking Through Capstone}
Your PA, client, and TAs are professional contacts who may be valuable throughout your career. After the project ends, they may serve as references, provide job leads, or connect you with their professional networks. Maintain the relationship:

\begin{itemize}[nosep,leftmargin=1.5em]
\item Send a thank-you email after the project closes.
\item Connect on LinkedIn with a personalized note (not the default ``I'd like to add you to my network'').
\item Provide occasional career updates (new job, graduate school acceptance, published paper).
\item \textbf{Always ask permission before listing someone as a professional reference.}
\end{itemize}

\section{LinkedIn Profile Optimization}
Your LinkedIn profile is often the first thing a recruiter sees. Capstone provides excellent material:

\textbf{Headline}: go beyond ``Student at Colorado School of Mines.'' Use: ``Mechanical Engineering Senior | Capstone: Solar-Powered Irrigation Control System | Seeking Full-Time Roles in Mechatronics.''

\textbf{Experience section}: add your capstone project as a position. Title: ``Capstone Design Engineer'' or ``Senior Design Team Lead.'' Organization: ``Colorado School of Mines; EDNS 491/492.'' Description: 3--5 bullet points quantifying your contributions, just like your resume.

\textbf{Projects section}: add the capstone project with a detailed description, your specific role, the technical skills used, and the outcome. Link to any public documentation (poster PDF, team video, GitHub repository).

\textbf{Skills section}: add the specific technical skills demonstrated in capstone: SolidWorks, KiCad, MATLAB, Python, FEA, PCB Design, PID Control, Project Management, Technical Writing, etc.

\textbf{Featured section}: upload your Showcase poster, a photo of your prototype, or a short video clip of the demonstration. Visual content dramatically increases profile engagement.

\section{Professional References}
When listing professional references on a job application:

\begin{itemize}[nosep,leftmargin=1.5em]
\item \textbf{Ask permission first}: never list someone as a reference without their explicit consent. Contact them before you submit the application, tell them about the position, and ask if they are willing to serve as a reference.
\item \textbf{Prepare your references}: send them your resume, a description of the position, and 2--3 key points you would like them to emphasize (technical skills, leadership, problem-solving ability). A reference who knows what to highlight is far more effective than one caught off guard.
\item \textbf{Best capstone references}: your PA (who observed your engineering and professional growth over two semesters), your client (who can speak to your professionalism and technical contributions), and your TA (who can speak to domain expertise). A reference from a capstone PA carries significant weight because PAs observe students in a professional context, not just an academic one.
\item \textbf{Thank your references}: after the application process, send a brief thank-you note regardless of the outcome. If you got the job, let them know.
\end{itemize}

\section{Technical Interviews and Capstone}
Many engineering employers conduct technical interviews that assess your depth of understanding. Your capstone project is excellent material for these conversations. Be prepared to discuss:

\textbf{System architecture}: draw a block diagram of your system on a whiteboard. Explain each subsystem, the interfaces between them, and why the architecture was organized this way. This demonstrates systems thinking.

\textbf{Key engineering decisions}: why did you choose this sensor over alternatives? What trade-offs did you make? What was the sensitivity analysis? This demonstrates engineering judgment, not just execution.

\textbf{Analysis methodology}: walk through a calculation from the Engineering Calculation Package. State the assumptions, show the approach, and explain the conclusion. The interviewer is assessing whether you understand the physics, not whether you memorized the formula.

\textbf{Test results and failures}: what did you test? What passed? What failed? How did you diagnose and resolve failures? Employers value engineers who can troubleshoot as much as engineers who can design. Discussing a failure you resolved intelligently is more impressive than claiming everything worked perfectly.

\textbf{Lessons learned}: what would you do differently? This question assesses self-awareness and growth mindset. A thoughtful answer (``I would have ordered the PCB two weeks earlier to leave margin for a board revision'') demonstrates maturity.

\section{Portfolio Development}
Create a visual portfolio documenting your capstone design process. Include: the problem statement, concept sketches, CAD renderings, fabrication photos, test setups, data visualizations, and the final product. A well-curated portfolio communicates your engineering process more effectively than a resume alone. Host it on a personal website, a GitHub repository with a README, or a LinkedIn featured section.


\section{Translating Capstone Experience to Different Industries}
Your capstone project is in a specific domain, but the skills transfer broadly:

\textbf{Project management}: ``Managed a 5-person interdisciplinary team through a 2-semester development cycle with 3 formal design reviews.'' This applies to any industry.

\textbf{Client communication}: ``Presented technical recommendations to a non-technical external client and incorporated feedback into 3 design revisions.'' Consulting, product management, and sales engineering all value this.

\textbf{Technical problem-solving}: ``Diagnosed a ground-loop-induced sensor error and implemented a star-grounding solution that reduced measurement noise by 85\%.'' This demonstrates systematic debugging; applicable to any technical role.

\textbf{Budget management}: ``Managed a \$2,500 project budget with 15\% contingency, completing the project 8\% under budget.'' Finance awareness is valued in every engineering role.

\textbf{Documentation}: ``Authored a 45-page Final Design Report including engineering analysis, test procedures, and an operations manual.'' Technical writing skills are consistently rated as the most important non-technical skill for new engineers.

\section{The Career Day Strategy}
Career Day is not a job fair where you passively hand out resumes. It is a networking event where you actively seek connections:

\textbf{Before Career Day}: research the companies attending. Identify 5--8 that match your interests. Review their job postings. Prepare specific questions for each (``I saw your posting for an embedded systems engineer; my capstone project involves ESP32-based control systems; could you tell me about the team structure for that role?'').

\textbf{At Career Day}: lead with your elevator pitch. Make eye contact. Hand over your resume while describing your capstone project. Ask about the company's engineering challenges, not just their open positions. Get business cards or LinkedIn contact information.

\textbf{After Career Day}: within 48 hours, send a brief follow-up email or LinkedIn connection request to every contact you made. Reference something specific from your conversation: ``It was great talking with you about your thermal management challenges. I would love to learn more about the team.'' Most students skip this step; the ones who follow up stand out.



\section{The Job Search Timeline for Graduating Seniors}
The job search does not begin at graduation; it begins in SDI:

\textbf{SDI (Fall semester)}: update resume with capstone project description (even early-stage). Attend Career Day. Apply to summer internships if applicable. Begin LinkedIn profile optimization.

\textbf{SDII (Spring semester)}: apply to full-time positions starting in January--February. Use the capstone elevator pitch in interviews. Request references from PA and client. Update resume with CDR and testing accomplishments. Attend spring career events.

\textbf{Post-Showcase}: update resume and LinkedIn with final project outcomes, quantified achievements, and Showcase photos. Send thank-you notes to references. Follow up with contacts from Career Day and networking events.

\textbf{The compound advantage}: students who start the job search in SDI (attending Career Day, networking, updating LinkedIn) have a significant advantage over students who start in Sprint~13 of SDII. Employers fill positions on rolling timelines; the best opportunities may close before graduation if you are not already in the pipeline.

\section{Common Resume Mistakes for Engineering Seniors}
\begin{enumerate}[nosep,leftmargin=1.5em]
\item \textbf{Listing duties instead of achievements}: ``Responsible for PCB design'' tells the employer nothing about your competence. ``Designed a 4-layer mixed-signal PCB with integrated power management, fabricated through JLCPCB, and completed board bring-up with zero rework'' tells a story of competence.
\item \textbf{No quantification}: ``Improved system performance'' vs. ``Reduced measurement error from $\pm$3.2\,\degC{} to $\pm$0.7\,\degC{}, a 78\% improvement.'' Numbers are more memorable and more credible than adjectives.
\item \textbf{Generic skills list}: ``Microsoft Office, teamwork, communication'' adds nothing. List specific, differentiating skills: ``SolidWorks (500+ hours), KiCad (PCB layout), Python (data analysis and automation), ANSYS (FEA/thermal), PID control system design and tuning.''
\item \textbf{Outdated formatting}: one page for undergraduates (two pages only if you have extensive research or work experience). Clean, consistent formatting. No photos, no decorative borders, no colored backgrounds. Let the content speak.
\item \textbf{Spelling and grammar errors}: these immediately signal carelessness. Proofread. Then have someone else proofread. Then proofread again.
\end{enumerate}



\section{The FE Exam: Your First Step to Licensure}
The Fundamentals of Engineering (FE) exam is a 5-hour-and-20-minute, computer-based exam administered by NCEES. Passing it makes you an Engineer Intern (EI) or Engineer-in-Training (EIT); the first step toward PE licensure (Chapter~22).

\textbf{When to take it}: during your senior year or within 1--2 years of graduation. The material is freshest while you are still in school. Many states allow you to register before graduation.

\textbf{What it covers}: mathematics, probability/statistics, engineering economics, ethics, statics, dynamics, thermodynamics, fluid mechanics, heat transfer, materials, and your discipline-specific topics (ME, CE, EE, etc.). The exam is open-reference (NCEES provides a digital reference handbook), but speed matters; you have $\sim$3 minutes per question.

\textbf{How to prepare}: use the NCEES Practice Exam (\$50, official questions), review the NCEES FE Reference Handbook (free download), and study weak areas 2--3 months before the exam. Campus review sessions may be available through your department.

\textbf{Why it matters for your career}: PE licensure is legally required for offering engineering services to the public, signing engineering documents, and holding certain government and consulting positions. Starting the licensure path now (by passing the FE) demonstrates professionalism and forward planning to employers.

\section{Building Your Professional Identity Beyond Capstone}
Capstone is a launching pad, not a destination. Continue developing your professional identity after graduation:

\textbf{Join a professional society}: ASME, IEEE, ASCE, AIChE, SAE, or your discipline's equivalent. Student membership is free or discounted. Professional membership provides: conferences, publications, networking, and professional development credits needed for PE maintenance.

\textbf{Seek mentorship}: identify 2--3 professionals (former PAs, clients, managers, senior engineers) who can advise your career decisions. A mentor who has navigated the path you are starting can save you years of trial and error.

\textbf{Document your work}: maintain a portfolio of your best engineering work (capstone project, research, personal projects). Update it with each significant project. A well-maintained portfolio differentiates you from candidates with identical resumes.

\textbf{Continue learning}: the half-life of engineering knowledge is $\sim$5 years. Technologies, standards, and best practices evolve continuously. Read technical publications, attend conferences, and pursue relevant certifications (PMP, Six Sigma, AWS, specific software certifications).



\section{Capstone as Interview Material: Specific Scenarios}
Prepare to discuss these specific capstone experiences in interviews:

\textbf{Technical decision under uncertainty}: ``We had to select between two sensor technologies before completing all our analysis. I gathered the available data, presented a comparison matrix to the team, and recommended Option~B based on its superior noise performance and lower integration risk. We validated this choice during testing, and the sensor met requirements with margin.''

\textbf{Conflict resolution}: ``Two team members disagreed on the motor controller architecture. I facilitated a structured discussion using our decision matrix framework, ensuring both perspectives were evaluated against the requirements. We reached consensus on a hybrid approach that incorporated the strengths of both proposals.''

\textbf{Schedule recovery}: ``Our PCB had a layout error that required a second revision, adding three weeks to our timeline. I reorganized the sprint backlog to parallelize firmware development with the board re-fabrication, so we recovered two of the three lost weeks and still completed testing on schedule.''

\textbf{Client management}: ``Our client requested a feature mid-project that would have required redesigning the power subsystem. I prepared a scope impact analysis showing the 4-week schedule extension and \$200 cost increase, and presented three options: add the feature and extend the schedule, defer it to future work, or implement a simplified version within the current timeline. The client chose the simplified version.''

\textbf{Technical failure}: ``Our first integration test failed because the motor's EMI was corrupting the sensor readings. I systematically diagnosed the issue using the signal walk technique; measuring at each stage of the signal chain until I localized the noise to the shared ground path. We implemented star grounding, which reduced the noise by 90\% and brought sensor accuracy within specification.''

Each of these follows the STAR framework (Chapter~20) and demonstrates a specific professional competency. Prepare 5--7 such stories before interview season.

\section{Networking Beyond Capstone: The Long Game}
Your professional network grows throughout your career. Capstone provides your initial nodes:

\textbf{Your PA}: a professor or industry professional who knows your engineering capabilities and work ethic from 8 months of close observation. This is one of the strongest references you will ever have. Maintain the relationship.

\textbf{Your client}: an industry professional who worked with you on a real project. Client references carry enormous weight with employers because they demonstrate that you can operate in a professional (not just academic) context.

\textbf{Your teammates}: your future professional network includes your classmates. The teammate who struggles through a tough CDR with you may be the hiring manager at your dream company in 10 years. Treat every professional relationship as long-term.

\textbf{Guest speakers and TAs}: capstone may include guest lectures from industry professionals. Introduce yourself, ask a thoughtful question, and follow up with a LinkedIn connection. These are low-effort, high-value networking opportunities.

\textbf{Alumni network}: Mines alumni are disproportionately supportive of fellow graduates. Leverage the Mines alumni network (LinkedIn, alumni events, department connections) when job searching. A cold outreach to a Mines alum is substantially more likely to get a response than a cold outreach to a stranger.



\section{Cover Letters That Reference Capstone}
When applying for positions where a cover letter is requested, capstone provides excellent material for the body paragraphs:

\textbf{Opening paragraph}: state the position you are applying for and how you learned about it. If a capstone connection (PA, client, TA, Career Day contact) referred you, mention them by name: ``Dr.~Smith, who served as my capstone project advisor, recommended I apply for this position.''

\textbf{Body paragraph 1 (technical alignment)}: connect your capstone technical skills to the job requirements. ``During my capstone project, I designed and implemented a PID-based temperature control system using an ESP32 microcontroller, which directly aligns with your team's work in embedded control for industrial automation.''

\textbf{Body paragraph 2 (professional skills)}: highlight soft skills demonstrated in capstone. ``As Scrum Master for a 5-person interdisciplinary team, I facilitated sprint planning, tracked deliverables, and managed communication with our external industry client across two semesters. This experience prepared me for the collaborative, client-facing work described in your posting.''

\textbf{Closing paragraph}: express enthusiasm and request an interview. Reference a specific aspect of the company that genuinely interests you (not generic flattery).

\section{Following Up After Applications and Interviews}
\textbf{After submitting an application}: if you have a contact at the company (from Career Day, LinkedIn, or a PA introduction), send a brief note: ``I recently applied for [position]. I wanted to let you know because of our conversation about [topic] at Career Day. I believe my capstone experience in [domain] aligns well with the role.''

\textbf{After an interview}: send a thank-you email within 24 hours to every interviewer. Reference a specific topic from the conversation: ``I enjoyed discussing your team's approach to thermal management; it reminded me of the challenges we solved in our capstone project's enclosure design.'' Keep it brief (3--5 sentences).

\textbf{After a rejection}: respond graciously. ``Thank you for letting me know. I appreciated the opportunity to learn about [company] and I would welcome the chance to be considered for future openings.'' Professional graciousness is memorable; bitterness is also memorable, but in the wrong way.

\textbf{After an offer}: before accepting, discuss the offer with a trusted advisor (PA, mentor, parent, career counselor). Negotiate respectfully if appropriate (salary, start date, relocation assistance). Once accepted, notify all other companies where you have active applications.



\section{Salary Negotiation for New Graduates}
When you receive a job offer, negotiation is expected and appropriate:

\textbf{Research market rates}: use Glassdoor, Levels.fyi, BLS Occupational Outlook Handbook, and Mines Career Center salary data to understand the range for your role, location, and experience level. Know the 25th, 50th, and 75th percentile salaries.

\textbf{Consider total compensation}: salary is one component. Also evaluate: signing bonus, annual bonus structure, stock/equity (for tech companies), retirement matching (401k match), health insurance quality and cost, paid time off, remote work flexibility, professional development budget (conference attendance, certifications), and relocation assistance. A \$75k offer with \$5k signing bonus, 6\% 401k match, and full remote flexibility may be more valuable than an \$82k offer with no bonus, 3\% match, and required 5-day in-office.

\textbf{How to negotiate}: express enthusiasm for the role first (``I'm very excited about this opportunity''). Then present your case with data: ``Based on my research, the market rate for this role in [city] is \$X--\$Y. Given my capstone experience in [relevant domain] and my [specific skill], I would like to discuss a salary of \$Z.'' Be specific, not vague. Be prepared for the employer to counter-offer.

\textbf{What is negotiable for new graduates}: salary (usually 5--10\% above the initial offer is achievable), start date, signing bonus, relocation assistance, and professional development budget. What is usually not negotiable: vacation days (often set by company policy for all new hires), job title (set by the HR classification system), and work location (set by the team's needs).

\textbf{When not to negotiate}: if the offer is already at or above market rate and includes strong benefits, aggressive negotiation may signal unrealistic expectations. Read the situation: a startup offering \$70k may have no room to move; a Fortune 500 offering \$70k almost certainly does.



\begin{figure}[htbp]
\centering
\begin{tikzpicture}[
    node distance=0.3cm,
    block/.style={rectangle, draw=MinesBlue, fill=LightBG, text width=2.5cm, minimum height=0.6cm, align=center, font=\tiny\sffamily\bfseries, rounded corners=2pt, line width=0.5pt},
    arrow/.style={-{Stealth[length=1.5mm]}, MinesBlue, line width=0.5pt}
]
\node[block] (grad) at (0,0) {Graduation\\{\tiny BS/MS}};
\node[block, right=0.8cm of grad] (ei) {Engineer Intern\\{\tiny Pass FE Exam}};
\node[block, right=0.8cm of ei] (exp) {4 Years Experience\\{\tiny Supervised Practice}};
\node[block, right=0.8cm of exp] (pe) {Professional Engineer\\{\tiny Pass PE Exam}};

\draw[arrow] (grad) -- (ei);
\draw[arrow] (ei) -- (exp);
\draw[arrow] (exp) -- (pe);
\end{tikzpicture}
\caption{The path to Professional Engineer (PE) licensure. Each step builds on the previous one. The FE exam is best taken during or shortly after your senior year while course material is fresh.}\label{fig:pe-path}
\end{figure}

The licensure pathway (Figure~\ref{fig:pe-path}) is straightforward but time-bounded. The FE exam covers material from your undergraduate courses; every semester you wait after graduation, the material fades. Students who pass the FE during their senior year or within 6 months of graduation have the highest pass rates. The 4-year experience requirement begins accumulating as soon as you start working under a licensed PE's supervision; your capstone PA may qualify if they hold a PE license.

\section{Your First 90 Days on the Job}
The transition from capstone to your first engineering position involves a mindset shift:

\textbf{You are no longer graded}: performance is evaluated through project outcomes, peer feedback, and manager observation over months and years, not through rubrics on individual assignments. Consistency matters more than occasional brilliance.

\textbf{You do not know the codebase/system/process}: every new engineer spends 3--6 months learning the organization's tools, processes, standards, and culture. Ask questions. Take notes. The colleagues who seem to know everything were in your position 2--5 years ago.

\textbf{Capstone skills transfer directly}: the sprint cadence, the stand-up meetings, the design reviews, the client communication, the documentation discipline; these are standard industry practices. You are better prepared than you think.


\section{Practice Questions}
\begin{enumerate}[leftmargin=1.5em]
\item Write three resume bullet points describing your capstone contributions (real or hypothetical). Each must include a quantified achievement.
\item Prepare a STAR response for: ``Tell me about a time you had to make a difficult technical decision under time pressure.'' Use a capstone scenario.
\item An interviewer asks ``What was the biggest failure in your capstone project?'' Write a response that demonstrates engineering maturity rather than defensiveness.
\end{enumerate}


